\par\medskip\noindent\textbf{Cómo leer la salida.}
{\small\setlength{\tabcolsep}{4pt}\renewcommand{\arraystretch}{1.15}
\begin{xltabular}{\linewidth}{@{}>{\raggedright\arraybackslash}p{0.24\linewidth}>{\raggedleft\arraybackslash}p{0.15\linewidth}>{\raggedright\arraybackslash}X@{}}
\toprule Indicador & Valor & Interpretación \\ \midrule \endfirsthead
\toprule Indicador & Valor & Interpretación \\ \midrule \endhead
\bottomrule \endlastfoot
\texttt{beta}, \texttt{s} & p. ej. FB 1,235; 1,053 & Parámetros del modelo de mercado de cada empresa. \\
FB: \texttt{AR(0)}, \texttt{t} & $-5{,}05$; $-4{,}74$ & Facebook cayó 5\,\% más de lo esperado el lunes posterior a la noticia. \\
FB: \texttt{CAR(0,+5)} & $-10{,}25$ ($t=-3{,}88$) & En una semana la pérdida anormal acumulada llegó a 10\,\%. \\
AAPL, AMZN, GOOG & $\lvert t\rvert<1{,}5$ & Ningún rendimiento anormal de las otras empresas es significativo. \\
Cartera de pares: \texttt{CAR(0,+5)}, $t$ & $-0{,}08$; $-0{,}04$ & Como grupo, las otras tecnológicas no sufrieron pérdidas: no hubo contagio. \\
Correlaciones de residuos & 0,30 a 0,47 & Los residuos de las empresas están correlacionados (shocks sectoriales comunes). \\
Prueba ingenua: ee & 1,547 & Error estándar que supone independencia entre empresas: es demasiado pequeño. \\
Cociente cartera/ingenuo & 1,30 & Ignorar la correlación subestima el error estándar en 30\,\%: con eventos en la misma fecha hay que usar la cartera. \\
FB menos pares: CAR, $t$, $p$ & $-10{,}23$; $-4{,}35$; $1{,}3\times10^{-5}$ & La pérdida fue específica de Facebook. \\
\end{xltabular}}
